\section{Challenge 4: Sắp xếp mảng một chiều}
\label{sec:challenge4}

\begin{requirementbox}{Yêu cầu bài toán}
Cho $N$ bản ghi $(v, label)$, chương trình cần sắp xếp ổn định theo $v$ giảm dần, sau đó theo \func{label} tăng dần từ điển. Nếu cả hai khóa bằng nhau, thứ tự xuất hiện ban đầu phải được giữ nguyên. Đề bài không cho phép sử dụng các hàm sắp xếp dựng sẵn của Python.
\end{requirementbox}

\subsection{Hiện thực đã xây dựng}
Trong \codepath{challenge/challenge4.py}, mỗi bản ghi được biến đổi thành \func{(-v, label)}. Cách biểu diễn này cho phép dùng thứ tự tăng dần tự nhiên của tuple để đạt được yêu cầu $v$ giảm dần và \func{label} tăng dần. Thuật toán sắp xếp cục bộ là bottom-up merge sort do nhóm tự hiện thực.

Khi số lượng bản ghi đủ lớn, chương trình chia mảng thành các chunk liên tiếp. Mỗi worker sắp xếp một chunk bằng \func{bottom\_up\_merge\_sort}. Sau đó, tiến trình chính lần lượt merge các chunk đã sắp xếp bằng hàm \func{merge\_two} cho đến khi thu được danh sách cuối cùng.

\subsection{Phân tích song song hóa}
Merge sort cho phép tách giai đoạn sắp xếp cục bộ trên các đoạn độc lập. Đây là phần được song song hóa trong mã nguồn hiện tại. Giai đoạn merge phụ thuộc vào kết quả của các chunk, nên được thực hiện sau khi toàn bộ worker hoàn tất.

Thiết kế này ưu tiên tính đơn giản và độ tin cậy. Việc để tiến trình chính merge kết quả giúp giảm số lần truyền dữ liệu qua lại giữa các process. Đổi lại, khi $N$ rất lớn, giai đoạn merge tuần tự có thể trở thành nút thắt hiệu năng.

\subsection{Thiết kế thuật toán song song}
\begin{lstlisting}[caption={Pseudocode song song cho Challenge 4}]
MAIN(path):
    records = read_records(path)
    keys = [(-value, label) for each record]
    sorted_keys = parallel_sort(keys)
    return restore_original_value(sorted_keys)

parallel_sort(records):
    if len(records) < threshold:
        return bottom_up_merge_sort(records)

    ranges = split_array_ranges(records)
    chunks = pool.map(sort_range, ranges)
    return merge_chunks(chunks)
\end{lstlisting}

\subsection{Lưu đồ thuật toán}
\begin{figure}[H]
\centering
\begin{tikzpicture}[node distance=0.8cm and 1.0cm, scale=0.82, transform shape]
\node[flowstart] (start) {Bắt đầu};
\node[flowprocess, below=of start] (read) {Đọc $N$ bản ghi $(v, label)$};
\node[flowprocess, below=of read] (key) {Chuyển khóa thành $(-v, label)$};
\node[flowdecision, below=of key] (small) {$N$ nhỏ hơn ngưỡng?};
\node[flowprocess, below left=1.1cm and 1.8cm of small] (seq) {Sắp xếp tuần tự bằng merge sort};
\node[flowprocess, below right=1.1cm and 1.8cm of small] (par) {Chia chunk và sort song song};
\node[flowmerge, below=2.7cm of small] (merge) {Merge các chunk theo thứ tự};
\node[flowstart, below=of merge] (end) {Khôi phục $v$ và trả kết quả};
\draw[flowarrow] (start) -- (read);
\draw[flowarrow] (read) -- (key);
\draw[flowarrow] (key) -- (small);
\draw[flowarrow] (small) -- node[left] {Có} (seq);
\draw[flowarrow] (small) -- node[right] {Không} (par);
\draw[flowarrow] (seq) |- (merge);
\draw[flowarrow] (par) |- (merge);
\draw[flowarrow] (merge) -- (end);
\end{tikzpicture}
\caption{Luồng xử lý song song của Challenge 4}
\label{fig:challenge4-flow}
\end{figure}

\subsection{Giải thích vai trò các bước}
Challenge 4 gồm hai pha chính: sắp xếp cục bộ và merge toàn cục. Trước khi chia việc, mỗi bản ghi $(v, label)$ được đổi thành khóa $(-v, label)$. Phép đổi dấu $v$ biến yêu cầu sắp xếp giảm dần theo $v$ thành sắp xếp tăng dần theo khóa thứ nhất; khóa \func{label} vẫn giữ thứ tự tăng dần từ điển.

Trong pha song song, tiến trình chính chia mảng thành các chunk liên tiếp. Mỗi worker chạy merge sort tự hiện thực trên một chunk và trả về chunk đã sắp xếp. Các worker không tương tác trực tiếp với nhau vì thứ tự bên trong mỗi chunk chỉ phụ thuộc vào dữ liệu của chunk đó. Sau đó, tiến trình chính merge các chunk theo thứ tự. Đây là bước quyết định tính đúng của toàn bộ danh sách, vì các chunk riêng lẻ đã có thứ tự nhưng chưa tạo thành một thứ tự toàn cục.

Vai trò của bước merge đặc biệt quan trọng đối với tính ổn định. Khi hai bản ghi có khóa bằng nhau, phần tử đến từ nửa trái được chọn trước. Do các chunk ban đầu giữ thứ tự xuất hiện trong input và được merge theo thứ tự từ trái sang phải, các bản ghi trùng cả $v$ và \func{label} vẫn giữ thứ tự tương đối ban đầu.

\subsection{Tiềm năng tăng tốc và ví dụ minh họa}
Sắp xếp có chi phí \bigO{N \log N}. Khi $N$ lớn, việc sắp xếp cục bộ các chunk song song có thể giảm đáng kể thời gian của pha sort. Tuy nhiên, speedup bị giới hạn bởi pha merge ở tiến trình chính và chi phí truyền các chunk đã sắp xếp từ worker về. Vì vậy, tiềm năng tăng tốc tốt nhất thường xuất hiện ở input lớn, còn input nhỏ có thể không có lợi do overhead process.

Ví dụ với các bản ghi \func{(5, QA), (5, QB), (-1, ID), (9, QA)}, sau khi đổi khóa ta có \func{(-5, QA)}, \func{(-5, QB)}, \func{(1, ID)}, \func{(-9, QA)}. Nếu chia thành hai chunk, worker 1 sắp \func{(5, QA), (5, QB)} và worker 2 sắp \func{(-1, ID), (9, QA)} thành \func{(9, QA), (-1, ID)}. Bước merge toàn cục tạo kết quả \func{(9, QA), (5, QA), (5, QB), (-1, ID)}. Cặp $(5, QA)$ và $(5, QB)$ minh họa tiêu chí phụ theo \func{label}; nếu hai bản ghi hoàn toàn trùng khóa, thứ tự ban đầu được giữ nhờ merge ổn định.

\subsection{Tính ổn định}
Tính ổn định được bảo đảm trong bước merge. Khi hai phần tử có khóa bằng nhau, hàm merge chọn phần tử ở nửa trái trước nhờ điều kiện so sánh \func{a <= b}. Vì các chunk được tạo theo thứ tự xuất hiện ban đầu và được merge theo đúng thứ tự đó, các bản ghi trùng khóa vẫn giữ được thứ tự tương đối ban đầu.

\subsection{Dữ liệu, phụ thuộc và chi phí}
\begin{itemize}
  \item \textbf{Phân phối dữ liệu}: mảng bản ghi được chia thành các đoạn liên tiếp; worker xử lý một đoạn.
  \item \textbf{Phụ thuộc dữ liệu}: các chunk độc lập trong giai đoạn sort, nhưng phụ thuộc khi merge.
  \item \textbf{Cân bằng tải}: chia theo số lượng bản ghi giúp tải giữa các worker tương đối đồng đều.
  \item \textbf{Chi phí giao tiếp}: đáng kể vì mỗi worker trả về một danh sách đã sắp xếp; merge cuối có độ phức tạp \bigO{N \log p}.
\end{itemize}

\subsection{Đánh giá hiệu năng}
\begin{table}[H]
\centering
\caption{Kết quả đo hiệu năng cho Challenge 4}
\resizebox{\linewidth}{!}{%
\begin{tabular}{@{}c c c c c c@{}}
\toprule
$N$ & Số worker & $T_{\mathrm{seq}}$ (s) & $T_{\mathrm{par}}$ (s) & Speedup & Efficiency \\
\midrule
$10^4$ & \placeholder{p} & \placeholder{tseq} & \placeholder{tpar} & \placeholder{speedup} & \placeholder{eff} \\
$10^5$ & \placeholder{p} & \placeholder{tseq} & \placeholder{tpar} & \placeholder{speedup} & \placeholder{eff} \\
$10^6$ & \placeholder{p} & \placeholder{tseq} & \placeholder{tpar} & \placeholder{speedup} & \placeholder{eff} \\
\bottomrule
\end{tabular}
}
\end{table}

\paragraph{Nhận xét.}
\placeholder{Nêu rõ ngưỡng N nhỏ có thể chạy tuần tự nhanh hơn do overhead. Với N lớn, sort chunk song song có lợi, nhưng merge tuần tự và chi phí truyền danh sách khiến speedup không tuyến tính.}
